% 原书物理页=印刷页43--51；3.2 全节（含 p48 四幅子图与 p49 一幅 A2 图；p51 末第三个习题的 iv) 跨页至 p52 顶部，一并译至 3.3 标题前）.
\section{抽象根系}\label{sec:ssl-3-2}
\begin{definition}\label{def:ssl-3-2-1}
设 $V$ 是一个有限维 $\mathbb{R}$-向量空间. $V$ 上的一个\newterm{ssl-abstract-root-system}{抽象根系}{abstract root system}（\newterm{ssl-reduced-root-system}{既约根系}{reduced root system}）是指两个集合 $\Phi\subset V$ 与 $\Phi^{\vee}\subset V^{*}$，连同它们之间的一个双射 $\alpha\mapsto\alpha^{\vee}$，使得：
\begin{enumerate}[label=\roman*)]
\item $\Phi$ 有限、在 $\mathbb{R}$ 上生成 $V$，且 $0\notin\Phi$；
\item 对一切 $\alpha\in\Phi$，有 $\alpha^{\vee}(\Phi)\subset\mathbb{Z}$；
\item 对一切 $\alpha\in\Phi$，有 $\alpha^{\vee}(\alpha)=2$ 且 $s_{\alpha}(\Phi)=\Phi$，其中令 $s_{\alpha}(v):=v-\alpha^{\vee}(v)\alpha$；
\item 对一切 $\alpha\in\Phi$，有 $\Phi\cap\mathbb{R}\alpha=\{\pm\alpha\}$.
\end{enumerate}
$V$ 的维数称为根系的\newterm{ssl-root-system-rank}{秩}{rank of a root system}，而 $\alpha^{\vee}$ 称为与根 $\alpha$ 相伴的\newterm{ssl-coroot}{对偶根}{coroot}.
\end{definition}
\begin{example*}
设 $\mathfrak g$ 是一个半单 $K$-Lie 代数（$K$ 为特征零的代数闭域），$\mathfrak h$ 是一个 Cartan 子代数，$\Phi\subset\mathfrak h^{*}_{\mathbb Q}\subset\mathfrak h^{*}$ 是根组成的集合，而 $\Phi^{\vee}=\{H_{\alpha},\ \alpha\in\Phi\}\subset\mathfrak h_{\mathbb Q}\subset\mathfrak h$. 那么由前一个命题，
\[\bigl(V:=\mathbb{R}\otimes_{\mathbb Q}\mathfrak h^{*}_{\mathbb Q},\ \Phi,\ \Phi^{\vee},\ \alpha\mapsto\alpha^{\vee}:=H_{\alpha}\bigr)\]
是一个根系，我们把它记作 $\Phi(\mathfrak g,\mathfrak h)$.

\textbf{线性情形：}$\mathfrak g=\mathfrak{sl}_n(K)$，$\mathfrak h=\{$对角矩阵$\}$. 令 $e_1,\cdots,e_n$ 为 $\mathbb{R}^n$ 的典则基，$e_1^{*},\cdots,e_n^{*}$ 为其对偶基. 于是 $(\mathfrak g,\mathfrak h)$ 的根系可以描述如下：
\[
V=\Bigl\{\sum_{i=1}^{n}x_i e_i\in\mathbb{R}^{n},\ \sum_{i=1}^{n}x_i=0\Bigr\}\subset\mathbb{R}^{n}
\quad\text{且}\quad
V^{*}=\Bigl\{\sum_{i=1}^{n}x_i e_i^{*}\in(\mathbb{R}^{n})^{*},\ \sum_{i}x_i=0\Bigr\}\subset(\mathbb{R}^{n})^{*},
\]
其上的对偶配对由 $\mathbb{R}^n$ 与 $(\mathbb{R}^n)^{*}$ 之间的配对诱导，并且
\[
\Phi=\{\alpha_{ij}=(e_i-e_j),\ 1\leqslant i\neq j\leqslant n\}
\quad\text{且}\quad
\Phi^{\vee}=\{\alpha_{ij}^{\vee}=(e_i^{*}-e_j^{*}),\ 1\leqslant i\neq j\leqslant n\},
\]
双射是显然的. 在这一描述中，$e_i^{*}$ 对应于对角矩阵 $E_{ii}$（它属于 $\mathfrak{gl}_n$），而 $e_i$ 对应于 $\mathfrak h$ 上由 $\sum_k\lambda_k E_{kk}\mapsto\lambda_{ii}$ 给出的线性形式.

\textbf{辛情形：}$\mathfrak g=\mathfrak{sp}_{2n}(K)$，$\mathfrak h=\{$对角矩阵$\}$. $\mathfrak{sp}_{2n}$ 的对角矩阵形如 $X=\sum_{i=1}^{n}x_i(E_{ii}-E_{i+n,i+n})$. 记 $e_i$ 为 $\mathfrak h$ 上由 $X\mapsto x_i$ 给出的线性形式，$V$ 为以 $e_1,\cdots,e_n$ 为基的向量空间. 可以验证 $(\mathfrak g,\mathfrak h)$ 的根系具有如下形状：
\[
\Phi=\{(\pm e_i\pm e_j),\ i\neq j,\ 2e_i\}
\quad\text{且}\quad
\Phi^{\vee}=\{(\pm e_i^{*}+\pm e_j^{*}),\ i\neq j,\ e_i^{*}\},
\]
双射是显然的. 在这一描述中，诸 $e_i^{*}$ 是诸 $e_i$ 的对偶基，因而对应于对角矩阵 $E_{ii}-E_{i+n,i+n}$.

\textbf{奇正交情形：}$\mathfrak g=\mathfrak{so}(\psi_{2n+1})$，$\mathfrak h=\{$对角矩阵$\}$. $\mathfrak{so}(\psi_{2n+1})$ 的对角矩阵也形如 $X=\sum_{i=1}^{n}x_i(E_{ii}-E_{i+n,i+n})$. 仍记 $e_i$ 为 $\mathfrak h$ 上由 $X\mapsto x_i$ 给出的线性形式，$V$ 为以 $e_1,\cdots,e_n$ 为基的向量空间. 可以验证 $(\mathfrak g,\mathfrak h)$ 的根系具有如下形状：
\[
\Phi=\{(\pm e_i\pm e_j),\ i\neq j,\ e_i\}
\quad\text{且}\quad
\Phi^{\vee}=\{(\pm e_i^{*}+\pm e_j^{*}),\ i\neq j,\ 2e_i^{*}\},
\]
双射是显然的. 在这一描述中，诸 $e_i^{*}$ 是诸 $e_i$ 的对偶基，因而对应于对角矩阵 $E_{ii}-E_{i+n,i+n}$.

\textbf{偶正交情形：}$\mathfrak g=\mathfrak{so}(\psi_{2n})$（$n>1$），$\mathfrak h=\{$对角矩阵$\}$. $\mathfrak{so}(\psi_{2n})$ 的对角矩阵总还是形如 $X=\sum_{i=1}^{n}x_i(E_{ii}-E_{i+n,i+n})$. 再记 $e_i$ 为 $\mathfrak h$ 上由 $X\mapsto x_i$ 给出的线性形式，$V$ 为以 $e_1,\cdots,e_n$ 为基的向量空间. 可以验证 $(\mathfrak g,\mathfrak h)$ 的根系具有如下形状：
\[
\Phi=\{(\pm e_i\pm e_j),\ i\neq j\}
\quad\text{且}\quad
\Phi^{\vee}=\{(\pm e_i^{*}+\pm e_j^{*}),\ i\neq j\},
\]
双射是显然的. 在这一描述中，诸 $e_i^{*}$ 是诸 $e_i$ 的对偶基，因而对应于对角矩阵 $E_{ii}-E_{i+n,i+n}$.
\end{example*}
\begin{numpara}[Weyl 群]\label{np:ssl-3-2-2}
注意，上面定义 iii) 中引入的 $V$ 的自同态 $s_{\alpha}$ 是轴为 $\operatorname{Ker}(\alpha^{\vee})$、方向为 $\mathbb{R}\cdot\alpha$ 的\newterm{ssl-reflection}{反射}{reflection}. 换言之，
\[
s_{\alpha}^{2}=\mathrm{id}_{V},\qquad
\operatorname{Ker}(s_{\alpha}+\mathrm{id}_{V})=\mathbb{R}\alpha
\qquad\text{且}\qquad
\operatorname{Ker}(s_{\alpha}-\mathrm{id}_{V})=\operatorname{Ker}(\alpha^{\vee}).
\]
\end{numpara}
\begin{definition*}
由诸 $s_{\alpha}$（$\alpha\in\Phi$）生成的子群 $W\subset\mathrm{GL}(V)$ 称为该根系的\newterm{ssl-weyl-group}{Weyl 群}{Weyl group}.
\end{definition*}
由根系公理 iii)，有 $W\cdot\Phi\subset\Phi$. 事实上，由于 $\Phi$ 生成 $V$，给出 $W$ 在 $\Phi$ 上作用的态射 $W\longrightarrow\mathfrak{S}_{\Phi}$ 是单射. 特别地，$W$ 有限. 于是我们知道如何在 $V$ 上构造在 $W$ 下不变的 Euclid 内积（参见第~\ref{np:ssl-2-1-8}~段的习题）. 对于这样一个内积 $\langle\;,\;\rangle$，Weyl 群 $W$ 是 Euclid 空间 $(V,\langle\;,\;\rangle)$ 的等距群的一个子群. 特别地，此时反射 $s_{\alpha}$ 就是以 $\alpha^{\perp}$ 为轴的正交反射，因而形如
\[s_{\alpha}(v)=v-\frac{2\langle v,\alpha\rangle}{\langle\alpha,\alpha\rangle}\,\alpha.\]
\begin{example*}
设 $\Phi$ 是 $(\mathfrak g,\mathfrak h)$ 的根系. 我们已经看到，Killing 型在 $\mathfrak h_{\mathbb R}$ 上诱导一个 Euclid 内积. 由对偶性可把它搬到 $\mathfrak h^{*}_{\mathbb R}$ 上. 由于
\[\alpha(H)=\frac{2}{\langle H_{\alpha},H_{\alpha}\rangle}\langle H_{\alpha},H\rangle,\]
我们得到内积
\[\langle\alpha,\beta\rangle:=\frac{4\langle H_{\alpha},H_{\beta}\rangle}{\langle H_{\alpha},H_{\alpha}\rangle\,\langle H_{\beta},H_{\beta}\rangle}.\]
于是
\[s_{\alpha}(\beta)=\beta-\beta(H_{\alpha})\alpha=\beta-\frac{\langle H_{\alpha},H_{\alpha}\rangle}{2}\langle\alpha,\beta\rangle\,\alpha=\beta-\frac{2\langle\alpha,\beta\rangle}{\langle\alpha,\alpha\rangle}\,\alpha,\]
这表明 $s_{\alpha}$ 是以 $\alpha^{\perp}$ 为轴的正交反射，从而 $\langle\;,\;\rangle$ 是 $W$-不变的.

\textbf{线性情形：}在 $\mathfrak{sl}_n(K)$ 的例子中，所得的内积是 $\mathbb{R}^n$ 上典则 Euclid 内积所诱导的内积的 $2n$ 倍\footnote{译注：原文如此. 按上式，$\mathfrak{sl}_n$ 情形在 $V=\mathfrak h^{*}_{\mathbb R}$ 上所得的内积其实是 $\mathbb R^n$ 典则 Euclid 内积所诱导者的 $1/(2n)$ 倍；恰为 $2n$ 倍的是 Killing 型在 $\mathfrak h_{\mathbb R}$ 上诱导的内积.}. 因此反射 $s_{\alpha_{ij}}$ 是 $\mathbb{R}^n$ 中以 $H_{ij}=\{(x_k)_k,\ x_i=x_j\}$ 为轴的正交反射 $s_{ij}$ 在 $V$ 上的限制. 于是 $s_{ij}(e_i)=e_j$，且对 $k\neq i,j$ 有 $s_{ij}(e_k)=e_k$. 由此把 $W$ 等同于通过置换典则基而作用的对称群 $\mathfrak{S}_n$.

\textbf{其他经典情形：}在 $\mathfrak{sp}_{2n}$、$\mathfrak{so}_{2n}\simeq\mathfrak{so}(\psi_{2n})$ 与 $\mathfrak{so}_{2n+1}\simeq\mathfrak{so}(\psi_{2n+1})$ 诸情形的每一个中，Killing 型都诱导``典则''内积（即使 $e_1,\cdots,e_n$ 成为标准正交基者）的一个倍数（分别为 $4n+4$、$4n-4$ 与 $4n-2$）. 对 $i\neq j$，与 $e_i-e_j$ 相伴的反射 $s_{ij}$ 交换 $e_i$ 与 $e_j$ 而固定其余的 $e_k$；与 $e_i+e_j$ 相伴的反射 $s'_{ij}$ 则交换 $e_i$ 与 $-e_j$ 而固定其余的 $e_k$.

在\textbf{辛情形与奇正交情形}中，与 $e_i$ 或 $2e_i$ 相伴的反射 $s_i$ 把 $e_i$ 送到 $-e_i$ 而保持其余的 $e_k$. 由此可把 $W$ 等同于半直积 $\mathfrak{S}_n\rtimes\{\pm1\}^n$，即每行、每列都恰有一个非零元素且该非零元素是一个正负号的那些矩阵构成的群.

在\textbf{偶正交情形}中，可以验证 $W$ 是 $\mathfrak{S}_n\rtimes\{\pm1\}^n$ 中由 $-1$ 的个数为偶数的矩阵组成的子群.
\end{example*}
\begin{numpara}[同构、直和、既约根系\footnote{译注：原文标题作 ``systèmes réduits''（既约根系）；本段实际讨论的是根系的不可约性，疑为 ``systèmes irréductibles'' 之误，照录原文.}]\label{np:ssl-3-2-3}
设 $\Phi=(\Phi,V,\Phi^{\vee},\alpha\mapsto\alpha^{\vee})$ 与 $\Phi'=(\Phi',V',\Phi'^{\vee},\alpha'\mapsto\alpha'^{\vee})$ 是两个根系.
\begin{itemize}
\item $\Phi$ 与 $\Phi'$ 之间的一个同构，是指一个 $\mathbb{R}$-线性同构 $f\colon V\xrightarrow{\ \sim\ }V'$，它在 $\Phi$ 与 $\Phi'$ 之间诱导一个双射，且使得对偶双射与诸映射 $\alpha\mapsto\alpha^{\vee}$ 相容.
\item 直和 $\Phi\oplus\Phi'$ 是指 $V\oplus V'$ 中的根系 $\Phi\sqcup\Phi'$，其对偶根为 $\Phi^{\vee}\sqcup\Phi'^{\vee}\subset V^{*}\oplus V'^{*}=(V\oplus V')^{*}$，双射是显然的. 特别地可见 $W(\Phi\oplus\Phi')=W(\Phi)\times W(\Phi')$.
\end{itemize}
如果一个根系不是两个根系的直和，就称它是\newterm{ssl-irreducible-root-system}{不可约}{irreducible root system}的.
\end{numpara}
\begin{exercise*}
设 $\Phi$ 是一个根系，并给 $V$ 配备一个 $W$-不变内积. 证明：$\Phi$ 可约当且仅当可以把 $\Phi$ 写成 $\Phi=\Phi_1\sqcup\Phi_2$，其中 $\Phi_1$ 与 $\Phi_2$ 正交.
\end{exercise*}
\begin{proposition*}
设 $\mathfrak g$ 是一个半单 $K$-Lie 代数.
\begin{enumerate}[label=\roman*)]
\item 根系 $\Phi(\mathfrak g,\mathfrak h)$ 在同构意义下不依赖于 Cartan 子代数 $\mathfrak h$ 的选取. 我们把它的同构类记作 $\Phi(\mathfrak g)$.
\item $\Phi(\mathfrak g_1\oplus\mathfrak g_2)=\Phi(\mathfrak g_1)\oplus\Phi(\mathfrak g_2)$. 此外，$\Phi(\mathfrak g)$ 不可约当且仅当 $\mathfrak g$ 是单的.
\end{enumerate}
\end{proposition*}
\begin{proof}
i) 我们已经看到，若 $\mathfrak h'$ 是另一个 Cartan 子代数，则存在 $g\in\operatorname{Int}(G)$ 使得 $\mathfrak h'=g\cdot\mathfrak h$. 于是 $g$ 诱导一个 $K$-线性同构 $\mathfrak h\xrightarrow{\ \sim\ }\mathfrak h'$. 对偶同构显然把 $\Phi(\mathfrak g,\mathfrak h')$ 映到 $\Phi(\mathfrak g,\mathfrak h)$ 上，而 $H_{\alpha}$ 的唯一性表明 $g$ 把 $\Phi(\mathfrak g,\mathfrak h)^{\vee}$ 保持双射地映到 $\Phi(\mathfrak g,\mathfrak h')^{\vee}$ 上. 总之，$g$ 诱导一个同构 $\Phi(\mathfrak g,\mathfrak h')\xrightarrow{\ \sim\ }\Phi(\mathfrak g,\mathfrak h)$.
ii) 第一个断言是显然的，它蕴涵：若 $\mathfrak g$ 不是单的，则 $\Phi(\mathfrak g)$ 不是不可约的. 反之，设 $\Phi(\mathfrak g,\mathfrak h)$ 可约. 那么特别地，$H_{\alpha}$（$\alpha\in\Phi$）所成的集合 $\Phi(\mathfrak g,\mathfrak h)^{\vee}$ 可写成两个子集 $\Phi^{\vee}_1\sqcup\Phi^{\vee}_2$ 的不交并，这两个子集对于 Killing 型在 $\mathfrak h$ 上的限制是正交的. 由此得到一个正交分解 $\mathfrak h=\mathfrak h_1\oplus\mathfrak h_2$，其中 $\mathfrak h_i=\operatorname{Vect}_{K}(\Phi^{\vee}_i)$. 令 $\mathfrak g_1:=\mathfrak h_1\oplus\sum_{\alpha\in\Phi_1}\mathfrak g_{\alpha}$，$\mathfrak g_2:=\mathfrak h_1\oplus\sum_{\alpha\in\Phi_2}\mathfrak g_{\alpha}$\footnote{译注：原文此处作 $\mathfrak h_1$；由上文分解 $\mathfrak h=\mathfrak h_1\oplus\mathfrak h_2$，应为 $\mathfrak h_2$.}. 于是 $\mathfrak g=\mathfrak g_1\oplus\mathfrak g_2$. 此外，对 $\alpha_1\in\Phi_1$ 与 $\alpha_2\in\Phi_2$，有 $\alpha_1+\alpha_2\notin\Phi$，故 $[\mathfrak g_{\alpha_1},\mathfrak g_{\alpha_2}]=0$. 于是 $[\mathfrak g_1,\mathfrak g_2]=0$，从而 $\mathfrak g_1$ 与 $\mathfrak g_2$ 都是理想，这表明 $\mathfrak g$ 不是单的.
\end{proof}
\begin{numpara}[对偶]\label{np:ssl-3-2-4}
在 $V$ 上固定一个 $W$-不变内积. 若借助 $\langle\;,\;\rangle$ 把 $V^{*}$ 与 $V$ 等同起来，则由公式
\[s_{\alpha}(v)=v-\frac{2\langle v,\alpha\rangle}{\langle\alpha,\alpha\rangle}\alpha=v-\alpha^{\vee}(v)\alpha\]
可推出 $\alpha^{\vee}=\dfrac{2}{\langle\alpha,\alpha\rangle}\alpha$（对称地，$\alpha=\dfrac{2}{\langle\alpha^{\vee},\alpha^{\vee}\rangle}\alpha^{\vee}$）. 于是我们看到 $s_{\alpha^{\vee}}=s_{\alpha}$ 保持 $\Phi^{\vee}$. 由此推出
\[(\Phi^{\vee},\ \Phi,\ \alpha^{\vee}\mapsto\alpha)\ \text{是}\ V^{*}\ \text{中一个根系}.\]
\end{numpara}
我们把它称为 $(\Phi,\Phi^{\vee},\alpha\mapsto\alpha^{\vee})$ 的\newterm{ssl-dual-root-system}{对偶根系}{dual root system}. 上面的讨论表明，映射 $s_{\alpha}\mapsto s_{\alpha^{\vee}}$ 诱导 Weyl 群之间的同构 $W(\Phi)\xrightarrow{\ \sim\ }W(\Phi^{\vee})$.
\begin{example*}
可以看出，$\mathfrak{sp}_{2n}$ 与 $\mathfrak{so}_{2n+1}$ 的根系互为对偶根系. 与此不同，$\mathfrak{sl}_n$ 与 $\mathfrak{so}_{2n}$ 的根系是``自对偶''的.
\end{example*}
\begin{numpara}[两条根之间的约束]\label{np:ssl-3-2-5}
取定两条不成比例且不正交的根 $\alpha,\beta\in\Phi$. 必要时把 $\alpha$ 换成 $-\alpha$ 并交换 $\alpha$ 与 $\beta$，我们可以假设
\[
\langle\alpha,\beta\rangle<0
\qquad\text{且}\qquad
\langle\alpha,\alpha\rangle\leqslant\langle\beta,\beta\rangle.
\]
这时有下列约束：
\[
0<\frac{\langle\alpha,\beta\rangle^{2}}{\langle\alpha,\alpha\rangle\langle\beta,\beta\rangle}<1
\qquad\text{其中}\qquad
\alpha^{\vee}(\beta)=\frac{2\langle\alpha,\beta\rangle}{\langle\alpha,\alpha\rangle}\in\mathbb{Z}
\qquad\text{且}\qquad
\beta^{\vee}(\alpha)=\frac{2\langle\alpha,\beta\rangle}{\langle\beta,\beta\rangle}\in\mathbb{Z},
\]
它们导出不等式
\[0<\alpha^{\vee}(\beta)\,\beta^{\vee}(\alpha)<4.\]
这样一来只剩下三种可能：
\begin{enumerate}[label=\roman*)]
\item $\alpha^{\vee}(\beta)=\beta^{\vee}(\alpha)=-1$，此时 $\langle\alpha,\alpha\rangle=\langle\beta,\beta\rangle$，且夹角 $(\alpha,\beta)=\dfrac{2\pi}{3}$；
\item $\alpha^{\vee}(\beta)=-2$，$\beta^{\vee}(\alpha)=-1$，此时 $2\langle\alpha,\alpha\rangle=\langle\beta,\beta\rangle$，且夹角 $(\alpha,\beta)=\dfrac{3\pi}{4}$；
\item $\alpha^{\vee}(\beta)=-3$，$\beta^{\vee}(\alpha)=-1$，此时 $3\langle\alpha,\alpha\rangle=\langle\beta,\beta\rangle$，且夹角 $(\alpha,\beta)=\dfrac{5\pi}{6}$.
\end{enumerate}
\end{numpara}
\begin{example*}
在 $\mathfrak{sl}_n$ 的根系中取 $\alpha=\alpha_{ij}$，$\beta=\alpha_{kl}$. 我们看到：
\begin{itemize}
\item $\{i,j\}\cap\{k,l\}=\varnothing\Rightarrow\alpha$ 与 $\beta$ 正交；
\item $|\{i,j\}\cap\{k,l\}|=1\Rightarrow$ 情形 i)（等长，且夹角为 $\dfrac{2\pi}{3}$）.
\end{itemize}
\end{example*}
\begin{example*}
对 $\mathfrak g=\mathfrak{sp}_{2n}$，取 $\alpha=e_i-e_j$ 与 $\beta=2e_j$，则处于情形 ii).
\end{example*}
\begin{exercise*}
验证：在偶正交情形，只要 $\langle\alpha,\beta\rangle<0$，就总处于情形 i). 与此相反，在奇正交情形，情形 ii) 可能出现.
\end{exercise*}
\begin{remark*}
只要 $\langle\alpha,\beta\rangle<0$，就有 $\alpha+\beta\in\Phi$. 事实上，可设 $\lVert\alpha\rVert\leqslant\lVert\beta\rVert$，而在上面各种情形中都有 $\alpha+\beta=s_{\beta}(\alpha)$.
\end{remark*}
\textbf{应用：秩 $2$ 根系的分类.}\quad 设 $\Phi\subset\mathbb{R}^2$ 是一个秩 $2$ 的根系，并假设 $\mathbb{R}^2$ 的典则内积是 $W$-不变的. 选取 $\alpha,\beta\in\Phi$ 不共线，使得 $\lVert\alpha\rVert\leqslant\lVert\beta\rVert$ 且使 $\dfrac{\langle\alpha,\beta\rangle}{\lVert\alpha\rVert\,\lVert\beta\rVert}$ 最小（于是 $\alpha$ 与 $\beta$ 的夹角在不为平角的前提下尽可能张开）. 此时有 $-1<\dfrac{\langle\alpha,\beta\rangle}{\lVert\alpha\rVert\,\lVert\beta\rVert}\leqslant 0$. 经过旋转与位似，可设 $\alpha=e_1$（典则基中的向量），必要时再作关于轴 $(e_1)$ 的对称，可设夹角 $(\alpha,\beta)$ 介于 $\dfrac{\pi}{2}$ 与 $\pi$ 之间. 这时有下列可能：
\begin{enumerate}[label=\roman*)]
\item 若 $\langle\alpha,\beta\rangle=0$，则 $\Phi=\{\pm\alpha,\pm\beta\}$. 必要时更换内积，可以看出 $\Phi$ 同构于下图中那个可约根系，它正是 $\mathfrak{sl}_2\oplus\mathfrak{sl}_2$ 的根系.
\begin{center}
\begin{tikzpicture}[>=Stealth]
\draw[dash pattern=on 0.5pt off 2.2pt, thin, <->] (-2.04,0) -- (2.04,0);
\draw[dash pattern=on 0.5pt off 2.2pt, thin, <->] (0,-2.04) -- (0,2.04);
\draw[thick, ->] (0,0) -- (1,0);
\draw[thick, ->] (0,0) -- (-1,0);
\draw[thick, ->] (0,0) -- (0,1);
\draw[thick, ->] (0,0) -- (0,-1);
\node at (-0.62,1.08) {$\beta=e_2$};
\node at (1.08,0.28) {$\alpha=e_1$};
\end{tikzpicture}
\end{center}
\end{enumerate}
若 $\langle\alpha,\beta\rangle<0$，则处于前面考虑过的三种情形之一.
\begin{enumerate}[label=\roman*),resume]
\item 当 $\alpha^{\vee}(\beta)=-1$ 时，$\beta$ 的极坐标为 $\bigl(1,\dfrac{2\pi}{3}\bigr)$. 我们注意到 $\alpha+\beta$ 也是根，于是 $\Phi\supset\{\pm\alpha,\pm\beta,\pm(\alpha+\beta)\}$. 而且，倘若还有别的根，就会与 $\dfrac{\langle\alpha,\beta\rangle}{\lVert\alpha\rVert\,\lVert\beta\rVert}$ 的最小性矛盾. 这一点在下图中更清楚. 所得的根系同构于 $\mathfrak{sl}_3$ 的根系.
\begin{center}
\begin{tikzpicture}[>=Stealth]
\draw[dash pattern=on 0.5pt off 2.2pt, thin, <->] (-2.04,0) -- (2.04,0);
\draw[dash pattern=on 0.5pt off 2.2pt, thin, <->] (0,-2.04) -- (0,2.04);
\draw[thick, ->] (0,0) -- (1,0);
\draw[thick, ->] (0,0) -- (-1,0);
\draw[thick, ->] (0,0) -- (60:1);
\draw[thick, ->] (0,0) -- (120:1);
\draw[thick, ->] (0,0) -- (240:1);
\draw[thick, ->] (0,0) -- (300:1);
\node at (1.08,0.24) {$\alpha$};
\node at (-0.79,0.93) {$\beta$};
\node at (1.00,1.13) {$\alpha+\beta$};
\end{tikzpicture}
\end{center}
\item 当 $\alpha^{\vee}(\beta)=-2$ 时，$\beta$ 的极坐标为 $\bigl(\sqrt{2},\dfrac{3\pi}{4}\bigr)$，于是 $\beta=e_2-e_1$. 利用反射 $s_{\alpha}$ 与 $s_{\beta}$，得到 $\Phi\supset\{\pm e_1,\pm e_2,\ \pm e_1\pm e_2\}$. 与上面一样，倘若还有别的根，它就会与 $\dfrac{\langle\alpha,\beta\rangle}{\lVert\alpha\rVert\,\lVert\beta\rVert}$ 的最小性矛盾. 所得的根系同构于 $\mathfrak{sp}_4$ 的根系，也同构于 $\mathfrak{so}_5$ 的根系.
\begin{center}
\begin{tikzpicture}[>=Stealth]
\draw[dash pattern=on 0.5pt off 2.2pt, thin, <->] (-2.04,0) -- (2.04,0);
\draw[dash pattern=on 0.5pt off 2.2pt, thin, <->] (0,-2.04) -- (0,2.04);
\draw[thick, ->] (0,0) -- (1,0);
\draw[thick, ->] (0,0) -- (-1,0);
\draw[thick, ->] (0,0) -- (0,1);
\draw[thick, ->] (0,0) -- (0,-1);
\draw[thick, ->] (0,0) -- (45:{sqrt(2)});
\draw[thick, ->] (0,0) -- (135:{sqrt(2)});
\draw[thick, ->] (0,0) -- (225:{sqrt(2)});
\draw[thick, ->] (0,0) -- (315:{sqrt(2)});
\node at (1.08,0.28) {$\alpha=e_1$};
\node at (-1.86,1.07) {$\beta=e_2-e_1$};
\node at (0.05,1.33) {$\beta+\alpha$};
\node at (1.56,1.33) {$\beta+2\alpha$};
\end{tikzpicture}
\end{center}
\item 当 $\alpha^{\vee}(\beta)=3$\footnote{译注：原文如此，未带负号；由前文情形 iii) 及 $\beta$ 的极坐标 $\bigl(\sqrt3,\frac{5\pi}{6}\bigr)$，此处应为 $\alpha^{\vee}(\beta)=-3$.}时，$\beta$ 的极坐标为 $\bigl(\sqrt{3},\dfrac{5\pi}{6}\bigr)$. 利用反射 $s_{\alpha}$ 与 $s_{\beta}$，得到 $\Phi\supset\{\pm\alpha,\ \pm\beta,\ \pm(\beta+\alpha),\ \pm(\beta+2\alpha),\ \pm(\beta+3\alpha),\ \pm(2\beta+3\alpha)\}$，并且与上面一样，可以验证这些就是全部的根. 所得的根系不是任何经典 Lie 代数的根系.
\begin{center}
\begin{tikzpicture}[>=Stealth]
\draw[dash pattern=on 0.5pt off 2.2pt, thin, <->] (-2.04,0) -- (2.04,0);
\draw[dash pattern=on 0.5pt off 2.2pt, thin, <->] (0,-2.04) -- (0,2.04);
\draw[thick, ->] (0,0) -- (1,0);
\draw[thick, ->] (0,0) -- (-1,0);
\draw[thick, ->] (0,0) -- (60:1);
\draw[thick, ->] (0,0) -- (120:1);
\draw[thick, ->] (0,0) -- (240:1);
\draw[thick, ->] (0,0) -- (300:1);
\draw[thick, ->] (0,0) -- (30:{sqrt(3)});
\draw[thick, ->] (0,0) -- (90:{sqrt(3)});
\draw[thick, ->] (0,0) -- (150:{sqrt(3)});
\draw[thick, ->] (0,0) -- (210:{sqrt(3)});
\draw[thick, ->] (0,0) -- (270:{sqrt(3)});
\draw[thick, ->] (0,0) -- (330:{sqrt(3)});
\node at (1.08,0.26) {$\alpha$};
\node at (-1.76,1.18) {$\beta$};
\node at (-0.62,1.18) {$\beta+\alpha$};
\node at (0.75,1.48) {$\beta+2\alpha$};
\node at (1.85,1.05) {$\beta+3\alpha$};
\node at (0,2.08) {$2\beta+3\alpha$};
\end{tikzpicture}
\end{center}
\end{enumerate}
\begin{numpara}[单根]\label{np:ssl-3-2-6}
取定一个不在任何根上取零值的线性形式 $f\in V^{*}$. 若 $f(\alpha)>0$，则称 $\alpha$ 是 $f$-正的（无歧义时简称\newterm{ssl-positive-root}{正根}{positive root}）. 我们把 $f$-正的根所成的集合记作 $\Phi^{+}$，并令 $\Phi^{-}:=-\Phi^{+}$. 这样得到一个分拆 $\Phi=\Phi^{+}\sqcup\Phi^{-}$（并不正交！！）. 若 $\alpha$ 是正的且不是两个正根之和，则称 $\alpha$ 是 $f$-单的（无歧义时简称\newterm{ssl-simple-root}{单根}{simple root}）. 我们把单根所成的集合记作 $\Delta=\Delta_f$.
\end{numpara}
\begin{example*}
取 $\mathfrak g=\mathfrak{sl}_n(K)$. 取 $\mathbb{R}^n$ 上的线性形式 $f\colon(x_i)_i\mapsto-\sum_i i\,x_i$ 并把它限制到 $V$ 上. 那么 $\Phi^{+}=\{\alpha_{ij},\ i<j\}$，而 $\Delta=\{\alpha_{i,i+1},\ i=1,\cdots,n-1\}$.
\end{example*}
\begin{theorem*}
采用上面的记号，有：
\begin{enumerate}[label=\roman*)]
\item $\Phi^{+}\subset\sum_{\alpha\in\Delta}\mathbb{N}\alpha$；
\item $\Delta$ 是 $V$ 的一组基.
\end{enumerate}
\end{theorem*}
\begin{proof}
i) 反证：取 $\beta\in\Phi^{+}\setminus\bigl(\sum_{\alpha\in\Delta}\mathbb{N}\alpha\bigr)$ 使 $f(\beta)$ 最小. 那么 $\beta\notin\Delta$，故 $\beta=\beta_1+\beta_2$，其中 $\beta_i\in\Phi^{+}$. 但此时 $f(\beta_i)<f(\beta)$，故由最小性 $\beta_i\in\sum_{\alpha\in\Delta}\mathbb{N}\alpha$，进而 $\beta\in\sum_{\alpha\in\Delta}\mathbb{N}\alpha$：矛盾.
ii) 由 i)，$\Delta$ 是生成组. 还须证明它是自由的. 取一个 $W$-不变内积，并首先注意到：对 $\Delta$ 中所有不同的 $\alpha,\beta$，有 $\langle\alpha,\beta\rangle\leqslant 0$. 事实上，若 $\langle\alpha,\beta\rangle>0$，我们在上一段中已指出 $\alpha-\beta$ 是根. 若它是正的，则 $\alpha=\beta+(\alpha-\beta)$ 与 $\alpha$ 的单性矛盾；否则与 $\beta$ 的单性矛盾.
现在，设 $\sum_{\alpha\in\Delta}\lambda_{\alpha}\alpha=0$ 是一个线性相关关系. 令
\[
\Delta^{+}=\{\alpha\in\Delta,\ \lambda_{\alpha}\geqslant 0\},
\qquad
\Delta^{-}=\Delta\setminus\Delta^{+}.
\]
把该关系写成
\[
\sum_{\alpha\in\Delta^{+}}\lambda_{\alpha}\alpha=-\sum_{\beta\in\Delta^{-}}\lambda_{\beta}\beta=:v
\]
的形式. 那么 $\langle v,v\rangle=\sum_{\alpha\in\Delta^{+},\,\beta\in\Delta^{-}}-\lambda_{\alpha}\lambda_{\beta}\langle\alpha,\beta\rangle\leqslant 0$，故 $v=0$. 于是 $\sum_{\alpha\in\Delta^{+}}\lambda_{\alpha}\alpha=0$，进而 $\sum_{\alpha\in\Delta^{+}}\lambda_{\alpha}f(\alpha)=0$，从而 $\lambda_{\alpha}=0,\ \forall\alpha\in\Delta^{+}$，同理 $\lambda_{\beta}=0,\ \forall\beta\in\Delta^{-}$.
\end{proof}
\begin{definition*}
$\Phi$ 的一组\newterm{ssl-root-base}{基}{base of a root system}是指像上面那样由某个 $f\in V^{*}$ 决定的 $f$-单根所成的集合 $\Delta=\Delta_f$.
\end{definition*}
人们自然会问：$f$ 与 $f'$ 满足什么条件时会有 $\Delta_f=\Delta_{f'}$？为此，回忆 $f$ 是 $X=V^{*}\setminus\bigcup_{\alpha}\operatorname{Ker}(\alpha)$ 中的元素，即 $V^{*}$ 中一组超平面之并（这些超平面称为\newterm{ssl-wall}{墙}{wall}）的补集中的元素. $X$ 的连通分支因此是一些开凸锥，称为\newterm{ssl-weyl-chamber}{Weyl 房}{Weyl chamber}，而 $f,f'$ 位于同一个 Weyl 房中当且仅当对每条根 $\alpha\in\Phi$，$f(\alpha)$ 与 $f'(\alpha)$ 同号. 由此得出：$\Delta_f=\Delta_{f'}$ 当且仅当 $f$ 与 $f'$ 位于同一个 Weyl 房中.
\begin{example*}
考虑前面见过的 $\mathfrak{sl}_3$ 的秩 $2$ 根系，并借助内积把 $V$ 与 $V^{*}$ 等同起来（从而把 $\operatorname{Ker}(\alpha)$ 与 $\alpha^{\perp}$ 等同起来）. 取 $f=\langle e_1+e_2,\,\cdot\,\rangle$，就得到基 $\Delta_f=\{\alpha,\beta\}$. 图中画出了墙以及与 $\Delta=\Delta_f$ 相伴的 Weyl 房 $C_{\Delta}$.
\begin{center}
\begin{tikzpicture}[>=Stealth]
\fill[gray!10] (0,0) -- (0,2.05) -- (30:2.05) -- cycle;
\draw[thin] (0,-2.05) -- (0,2.05);
\draw[thin] (30:2.05) -- (210:2.05);
\draw[thin] (150:2.05) -- (330:2.05);
\draw[thick, ->] (0,0) -- (1,0);
\draw[thick, ->] (0,0) -- (-1,0);
\draw[thick, ->] (0,0) -- (60:1);
\draw[thick, ->] (0,0) -- (120:1);
\draw[thick, ->] (0,0) -- (240:1);
\draw[thick, ->] (0,0) -- (300:1);
\node at (1.08,0.24) {$\alpha$};
\node at (-0.79,0.93) {$\beta$};
\node at (0.78,1.42) {$C_{\Delta}$};
\node at (240:2.05) [below right] {$\Delta=\{\alpha,\beta\}$};
\end{tikzpicture}
\end{center}
\end{example*}
在这个例子上可以观察到两个有趣的事实.
\begin{itemize}
\item 一方面，$W$ 由 $s_{\alpha}$ 与 $s_{\beta}$ 生成. 事实上，我们已经指出 $W=\mathfrak{S}_3$，而众所周知这个对称群由对换 $s_{\alpha}=(1,2)$ 与 $s_{\beta}=(2,3)$ 生成.
\item 另一方面，Weyl 房共有 $6=3!=|W|$ 个. 事实上，$W$ 在 $V^{*}$ 上的作用单传递地置换诸 Weyl 房. 等价地说，$W$ 在 $\Phi$ 上的作用单传递地置换 $\Phi$ 的诸基.
\end{itemize}
我们将看到，这两个事实在一般情形下仍然成立.
\begin{proposition}\label{prop:ssl-3-2-7}
设 $\Delta$ 是 $\Phi$ 的一组基. 把满足 $\alpha\in\Delta$ 的反射 $s_{\alpha}$ 称为\newterm{ssl-simple-reflection}{单反射}{simple reflection}.
\begin{enumerate}[label=\roman*)]
\item 每条根都是某条单根在若干单反射的乘积下的像.
\item Weyl 群由诸单反射生成.
\end{enumerate}
\end{proposition}
下面的引理将被多次使用.
\begin{lemma*}
设 $\alpha\in\Delta$ 是一条单根. 则 $s_{\alpha}\bigl(\Phi^{+}\setminus\{\alpha\}\bigr)=\Phi^{+}\setminus\{\alpha\}$.
\end{lemma*}
\begin{proof}
设 $\gamma$ 是异于 $\alpha$ 的正根. 在基 $\Delta$ 中把它写成 $\gamma=\sum_{\beta\in\Delta}n_{\beta}\beta$，其中 $n_{\beta}\in\mathbb{N}$. 我们有 $s_{\alpha}(\gamma)=\sum_{\beta\in\Delta}n'_{\beta}\beta$，其中当 $\beta\neq\alpha$ 时 $n'_{\beta}=n_{\beta}$，而 $n'_{\alpha}=n_{\alpha}-\alpha^{\vee}(\gamma)$. 但由于 $\gamma\neq\alpha$，存在 $\beta$ 使得 $n'_{\beta}=n_{\beta}>0$. 由定理的 i)，推出 $s_{\alpha}(\gamma)$ 不是\newterm{ssl-negative-root}{负根}{negative root}，因而是正的. 另外，显然有 $s_{\alpha}(\gamma)\neq\alpha$.
\end{proof}
\begin{proof}[命题~\ref{prop:ssl-3-2-7} 的证明]
i) 只须对正的 $\gamma$ 来做，因为若 $\gamma=w(\beta)$，则 $-\gamma=ws_{\beta}(\beta)$. 设 $\gamma\in\Phi^{+}$. 把 $\gamma$ 在基 $\Delta$ 中坐标之和 $h(\gamma)=\sum_{\beta}n_{\beta}$ 定义为 $\gamma$ 的\newterm{ssl-height}{高}{height of a root}，并对 $h(\gamma)$ 作归纳. 若 $h(\gamma)=1$，则 $\gamma$ 是单的，无须证明. 否则，可以找到 $\alpha\in\Delta$ 使 $\alpha^{\vee}(\gamma)>0$. 事实上，对 $W$-不变内积 $\langle\;,\;\rangle$ 有 $\alpha^{\vee}(\gamma)=\dfrac{2\langle\gamma,\alpha\rangle}{\langle\alpha,\alpha\rangle}$ 且 $\lVert\gamma\rVert^{2}=\sum_{\beta\in\Delta}n_{\beta}\langle\gamma,\beta\rangle>0$. 于是，与上面的引理一样，根 $s_{\alpha}(\gamma)$ 是正的，且 $h(s_{\alpha}(\gamma))=h(\gamma)-\alpha^{\vee}(\gamma)<h(\gamma)$. 因此可对高作归纳而得出结论.
ii) 按定义，$W$ 由所有反射 $s_{\gamma}$（$\gamma\in\Phi$）生成. 由 i)，可写 $\gamma=w(\alpha)$，其中 $w$ 是一些单反射的乘积而 $\alpha\in\Delta$. 由于对 $W$-不变内积而言 $s_{\gamma}$ 是以 $\gamma^{\perp}$ 为轴的正交反射，此时有 $s_{\gamma}=ws_{\alpha}w^{-1}$. 所以 $s_{\gamma}$ 确是单反射的乘积.
\end{proof}
\begin{theorem}\label{thm:ssl-3-2-8}
$W$ 在 $\Phi$ 上的作用单传递地置换 $\Phi$ 的诸基.
\end{theorem}
\begin{proof}
首先注意，若 $\Delta=\Delta_f$，则集合 $w(\Delta)=\Delta_{w^{*}(f)}$ 仍是 $\Phi$ 的基. 因此 $W$ 的作用确实置换 $\Phi$ 的诸基.
i) 先证该作用是传递的. 设 $\Delta'$ 是 $\Phi$ 的另一组基. 我们将对基数 $d(\Delta')=\bigl|\Phi^{+}(\Delta)\cap\Phi^{-}(\Delta')\bigr|$ 作归纳来证明 $\Delta'\in W\cdot\Delta$. 若 $d(\Delta')=0$，则 $\Delta=\Delta'$，断言显然. 故设 $d(\Delta')>0$，即 $\Phi^{+}(\Delta)\cap\Phi^{-}(\Delta')\neq\varnothing$. 那么也有 $\Phi^{-}(\Delta)\cap\Delta'\neq\varnothing$（因为否则将有 $\Delta'\subset\Phi^{+}(\Delta)$，从而 $\Phi^{+}(\Delta')=\Phi^{+}(\Delta)$，最终得 $\Delta'=\Delta$）. 于是可取 $\alpha'\in\Phi^{-}(\Delta)\cap\Delta'$. 由上面最后一个引理，有 $s_{\alpha'}\bigl(\Phi^{-}(\Delta')\bigr)=\Phi^{-}(\Delta')\setminus\{-\alpha'\}\cup\{\alpha'\}$. 又有 $s_{\alpha'}\bigl(\Phi^{-}(\Delta')\bigr)=\Phi^{-}\bigl(s_{\alpha'}(\Delta')\bigr)$，故 $\Phi^{+}(\Delta)\cap\Phi^{-}\bigl(s_{\alpha'}(\Delta')\bigr)=\Phi^{+}(\Delta)\cap\Phi^{-}(\Delta')\setminus\{-\alpha'\}$，最终 $d\bigl(s_{\alpha'}(\Delta')\bigr)=d(\Delta')-1$. 由归纳假设，存在 $w\in W$ 使 $s_{\alpha'}(\Delta')=w(\Delta)$，这就给出 $\Delta'=s_{\alpha'}w(\Delta)$，如所欲证.
ii) 现设 $w\in W$，$w\neq 1$. 把它写成单反射的一个乘积 $s_1s_2\cdots s_l$，且 $l$ 最小. 记 $\alpha_i\in\Delta$ 为使 $s_i=s_{\alpha_i}$ 的单根. 我们将证明 $w(\alpha_l)\in\Phi^{-}$，由此推出 $w(\Delta)\neq\Delta$.
反证，设 $w(\alpha_l)\in\Phi^{+}$. 记 $j$ 为使 $\beta_i:=s_is_{i+1}\cdots s_l(\alpha_l)\in\Phi^{+}$ 的最大整数 $i$. 由于 $s_l(\alpha_l)\in\Phi^{-}$，有 $j<l$，故 $\beta_{j+1}$ 有定义且是负的. 由于 $\beta_{j+1}=s_j(\beta_j)$ 而 $\beta_j\in\Phi^{+}$，上面最后一个引理保证 $\beta_j=\alpha_j$. 由此得出 $s_{\beta_j}=(s_j\cdots s_l)\,s_{\alpha_l}\,(s_j\cdots s_l)^{-1}$，把它改写为 $s_{j+1}\cdots s_l=s_j\cdots s_{l-1}$. 但这样一来就得到 $w=s_1\cdots s_{j-1}s_{j+1}\cdots s_{l-1}$，这是一个长 $l-2$ 的表达式，与 $l$ 的最小性矛盾.
\end{proof}
\begin{exercise*}
验证：$\{(e_i-e_{i+1}),\ i=1,\cdots,n-1\}\cup\{e_n\}$ 是 $\mathfrak{so}_{2n+1}$ 的根系的一组基，而 $\{(e_i-e_{i+1}),\ i=1,\cdots,n-1\}\cup\{2e_n\}$ 是 $\mathfrak{sp}_{2n}$ 的根系的一组基. 试找出 $\mathfrak{so}_{2n}$ 的根系的一组基.
\end{exercise*}
\begin{exercise*}
设 $\Phi$ 是一个根系，$\Delta$ 是 $\Phi$ 的一组基，并给定 $\Delta$ 的一个分拆 $\Delta=\Delta_1\sqcup\Delta_2$.
\begin{enumerate}[label=\roman*)]
\item 令 $V_i:=\operatorname{Vect}_{\mathbb R}(\Delta_i)$，$\Phi_i:=\Phi\cap V_i$. 证明 $\Phi_i$ 是 $V_i$ 中的一个根系（补全各对象），其 Weyl 群 $W_i$ 是由诸 $s_{\alpha}$（$\alpha\in\Delta_i$）生成的群.
\item 以下设 $\Delta_1\perp\Delta_2$（对某个 $W$-不变内积）. 证明 $W=W_1\times W_2$，且 $W_2$ 平凡地作用于 $\Delta_1$.
\item 仍在假设 $\Delta_1\perp\Delta_2$ 下，证明 $\Phi=\Phi_1\sqcup\Phi_2$〔可利用上面命题~\ref{prop:ssl-3-2-7} 的 i)〕，并由此得出 $\Phi$ 是根系 $\Phi_1$ 与 $\Phi_2$ 的直和.
\end{enumerate}
\end{exercise*}
\begin{exercise*}
设 $\Phi$ 是一个根系，$\Delta$ 是 $\Phi$ 的一组基，并设在 $V$ 上固定了一个 $W$-不变内积. 在 $\Phi$ 上定义一个偏序关系：规定 $\gamma\preceq\gamma'$ 当且仅当 $\gamma'-\gamma\in\sum_{\alpha\in\Delta}\mathbb{N}\alpha$. 设 $\beta$ 是该偏序下的一条极大根.
\begin{enumerate}[label=\roman*)]
\item 验证 $\beta$ 是正的，并证明对一切 $\alpha\in\Delta$ 有 $\langle\beta,\alpha\rangle\geqslant 0$.
\item 写 $\beta=\sum_{\alpha\in\Delta}n_{\alpha}\alpha$，令 $\Delta_1:=\{\alpha,\ n_{\alpha}=0\}$，$\Delta_2=\Delta\setminus\Delta_1$. 证明 $\Delta_1\perp\Delta_2$.
\item 证明：若 $\langle\gamma,\gamma'\rangle>0$，则 $\gamma\preceq\gamma'$ 或 $\gamma'\preceq\gamma$.〔先指出在这种情形下 $\gamma-\gamma'\in\Phi$.〕
\item 以下设 $\Phi$ 不可约，于是由前一习题，$\Delta$ 不是两个正交的真子集之并. 证明 $\beta$ 是 $\Phi$ 中唯一的极大元（称之为最高根）.〔若 $\beta'$ 是另一个极大元，利用 i) 与 ii) 证明 $\langle\beta',\beta\rangle>0$，再利用 iii) 得出 $\beta=\beta'$.〕
\end{enumerate}
\end{exercise*}
